% METHODS-ONLY EXTRACT of per_em_W1_wiley.tex — 2026-10-11 — for clean-room replication.
% The replicator may read THIS FILE and the protocol's DGM table only. Not the project's R/, results/, or the rest of the paper.
% If something needed to implement the estimator is missing here, that is a finding about the Methods section.

\section{Methods}\label{sec:methods}

An effect modifier is a factor for which the treatment effect varies across subgroups. Formally, let the conditional average treatment effect (CATE) be denoted $\tau(x) = E[Y(1) - Y(0) \mid X = x]$, where $X$ is the effect modifier. When this function is non-constant, the average treatment effect observed in region $r$ depends on the distribution $F_r$ of the effect modifier in that region:
\begin{equation}
\bar{\tau}_r = \int \tau(x) \, dF_r(x).
\label{eq:regional_ate}
\end{equation}
Assessing pooling suitability therefore requires a measure of distributional difference between regions.

\subsection{Existing Approaches and Their Limitations}\label{sec:existing}

The most widely used measure for comparing covariate distributions across groups is the standardized mean difference (SMD), defined as
\begin{equation}
\text{SMD} = \frac{\bar{X}_1 - \bar{X}_2}{S_{\text{pooled}}},
\label{eq:smd}
\end{equation}
where $\bar{X}_r$ denotes the sample mean in region $r$ and $S_{\text{pooled}}$ is the pooled standard deviation.\cite{austin2011} Because treatment effects in clinical trials are typically summarized as mean differences, SMD provides a scale-free comparison directly relevant to the quantities that drive clinical decision-making. However, by construction SMD captures only differences in location. Two distributions with identical means but markedly different variances---for example, $N(50, 5^2)$ and $N(50, 15^2)$---yield $\text{SMD} = 0$ despite a threefold difference in spread.

Distributional distance measures address this limitation by comparing distributions beyond summary statistics. Empirical distribution function (EDF) tests such as the Kolmogorov--Smirnov (KS) statistic, $D_{\text{KS}} = \sup_x |F_1(x) - F_2(x)|$, and related measures (Cram\'{e}r--von Mises, Anderson--Darling) quantify differences in the full shape of distributions by comparing CDFs directly. Information-theoretic divergences such as Kullback--Leibler (KL) divergence,
\begin{equation}
D_{\text{KL}}(P \| Q) = \int p(x) \log \frac{p(x)}{q(x)} \, dx,
\label{eq:kl}
\end{equation}
take a density-based approach, measuring how one probability distribution differs from another.

However, each of these approaches has limitations when applied to comparing effect modifier distributions across regions. SMD captures only location differences and is blind to variance and shape. EDF-based statistics such as the KS statistic capture full distributional differences but lack a clinically interpretable scale, and the bound on treatment effect heterogeneity they admit carries a constant that clinical sources do not report (Section~\ref{sec:nabcd_metric}). KL divergence is asymmetric---comparing region A to region B yields a different value than comparing B to A---can diverge to infinity when empirical distributions do not fully overlap, and requires density estimation that is impractical with small regional samples. None of these measures offers a theoretical link that would translate a distributional distance into a bound on regional treatment effect differences whose constant is a per-unit clinical slope, the kind of quantity that treatment-by-covariate interaction analyses estimate. These gaps motivate a measure that (i) captures distributional features beyond location, (ii) provides a clinically interpretable scale, and (iii) offers such a link. The Wasserstein-1 distance and its clinical calibration, developed in the following subsections, satisfy all three requirements.

\subsection{Per-Effect-Modifier Wasserstein-1 Distance}\label{sec:nabcd_metric}

Among distributional distances that could capture the full distributional differences relevant to regional pooling, we adopt the Wasserstein-1 distance ($W_1$) for the clinical specifiability of the constant in its heterogeneity bound (Proposition~\ref{prop:heterogeneity} below). The $W_1$ distance (sometimes called the Earth Mover's Distance) admits three equivalent representations on the real line, each highlighting a different facet of its interpretation.\cite{villani2009,vallender1974}

Let $F_1$ and $F_2$ denote the cumulative distribution functions (CDFs) of an effect modifier in two regions. Then
\begin{equation}
W_1(F_1, F_2) = \int_{-\infty}^{\infty} |F_1(x) - F_2(x)| \, dx = \int_0^1 |F_1^{-1}(u) - F_2^{-1}(u)| \, du = \inf_{\pi \in \Pi(F_1, F_2)} \int |x - y| \, d\pi(x, y),
\label{eq:wasserstein}
\end{equation}
where $F_r^{-1}$ is the quantile function and $\Pi(F_1, F_2)$ is the set of probability couplings with marginals $F_1$ and $F_2$. The CDF form expresses $W_1$ as the area between the two CDFs (Figure~\ref{fig:nabcd_definition}); the quantile form as the integrated absolute difference of quantile functions; and the Kantorovich form as the minimum expected absolute displacement to transport one distribution onto the other. The three forms are mathematically equivalent.\cite{vallender1974} Unlike the standardized mean difference, $W_1$ responds to differences in variance, skewness, and other distributional features.\cite{panaretos2019}

\begin{figure}[t]
\centerline{\includegraphics[width=0.8\textwidth]{fig1_w1_definition.pdf}}
\caption{The Wasserstein-1 distance as the area between cumulative distribution functions, illustrated with a right-skewed Region~1 (gamma distribution) and a symmetric Region~2 (normal distribution). Panel~(A): probability densities. Panel~(B): cumulative distribution functions; the shaded area equals $W_1(F_1, F_2)$.\label{fig:nabcd_definition}}
\end{figure}

A property of $W_1$ central to our framework is that it carries the units of the underlying effect modifier. In each of the three equivalent forms in equation~(\ref{eq:wasserstein}), $W_1$ carries the units of $X$: the CDF form integrates a dimensionless difference of probabilities over $x$, the quantile form integrates differences of quantiles, and the coupling form integrates distances $|x - y|$. Consequently, $W_1(F_1, F_2)$ is expressed on the original scale of the effect modifier. This dimensional alignment supports direct clinical interpretation: a regional $W_1 = 5$~mmHg in systolic blood pressure means that optimally transporting one distribution onto the other moves probability mass an average distance of 5~mmHg. The same logic precludes a meaningful direct comparison of $W_1$ values across effect modifiers measured on different scales, and we therefore restrict attention in this paper to per-effect-modifier assessment (Section~\ref{sec:discussion} elaborates on this scope choice).

$W_1$ is a proper metric on the space of probability distributions with finite first moment, satisfying non-negativity, symmetry, identity of indiscernibles, and the triangle inequality.\cite{villani2009} It is also affine-equivariant in the variable units, satisfying $W_1(F_{aX+b}, F_{aY+b}) = |a| \cdot W_1(F_X, F_Y)$ for any $a \neq 0$. Beyond these basic properties, $W_1$ admits a Kantorovich--Rubinstein dual representation\cite{villani2009}
\begin{equation}
W_1(F_1, F_2) = \sup_{\|f\|_{\text{Lip}} \leq 1} \left|\int f \, dF_1 - \int f \, dF_2\right|,
\label{eq:dual}
\end{equation}
which expresses $W_1$ as the largest possible difference in expected values of a 1-Lipschitz function $f$ between the two distributions (the notation $\|f\|_{\text{Lip}} \leq 1$ means $|f(x) - f(y)| \leq |x - y|$ for all $x, y$). If the conditional average treatment effect function $\tau(x)$ is Lipschitz continuous with constant $L_{\text{clinical}}$, so that $|\tau(x) - \tau(y)| \leq L_{\text{clinical}} |x - y|$ for all $x, y$, then $\tau / L_{\text{clinical}}$ is 1-Lipschitz and applying equation~(\ref{eq:dual}) yields the central result of this section:

\begin{proposition}\label{prop:heterogeneity}
\textup{(Connection to heterogeneity).} Suppose the CATE function $\tau$ is the same in every region and depends on the effect modifier $X$ alone. If $\tau$ is Lipschitz continuous with constant $L_{\text{clinical}}$, then for any two regions with effect modifier distributions $F_1$ and $F_2$,
\begin{equation}
|\bar{\tau}_1 - \bar{\tau}_2| \leq L_{\text{clinical}} \cdot W_1(F_1, F_2).
\label{eq:heterogeneity_bound}
\end{equation}
\end{proposition}

\begin{proof}
Apply equation~(\ref{eq:dual}) with $f = \tau / L_{\text{clinical}}$, which is 1-Lipschitz by assumption. This yields $\left|\int (\tau / L_{\text{clinical}}) \, d(F_1 - F_2)\right| \leq W_1(F_1, F_2)$. Combining with equation~(\ref{eq:regional_ate}) for $\bar{\tau}_r$ and rearranging yields equation~(\ref{eq:heterogeneity_bound}).
\end{proof}

\noindent Equation~(\ref{eq:heterogeneity_bound}) is not the only bound of its form. Integration by parts gives $\bar{\tau}_1 - \bar{\tau}_2 = -\int (F_1(x) - F_2(x)) \, d\tau(x)$, so for a CATE of bounded variation on a bounded support, $|\bar{\tau}_1 - \bar{\tau}_2| \leq \mathrm{TV}(\tau) \cdot D_{\text{KS}}$ as well, and $W_1 \leq W_2$ makes $L_{\text{clinical}} \cdot W_2$ a valid but looser bound.\cite{villani2009} What distinguishes equation~(\ref{eq:heterogeneity_bound}) is its constant. $L_{\text{clinical}}$ is a change in treatment effect per unit of the effect modifier, the kind of quantity that treatment-by-covariate interaction analyses estimate; $\mathrm{TV}(\tau)$ is the total excursion of the CATE over the support, which reduces to a clinically stated quantity only for a monotone CATE. Where both constants are available, neither bound dominates in general. On a support $[a, b]$, $W_1 \leq (b - a)\,D_{\text{KS}}$ and $\mathrm{TV}(\tau) \leq L_{\text{clinical}}\,(b - a)$, and $L_{\text{clinical}} \cdot W_1 \leq \mathrm{TV}(\tau) \cdot D_{\text{KS}}$ holds exactly when $W_1 / D_{\text{KS}} \leq \mathrm{TV}(\tau) / L_{\text{clinical}}$. For a CATE whose slope has constant magnitude $L_{\text{clinical}}$, such as a linear CATE, $\mathrm{TV}(\tau) = L_{\text{clinical}}(b - a)$ and equation~(\ref{eq:heterogeneity_bound}) is never the looser of the two; the Kolmogorov--Smirnov bound gains as the CATE's slope falls below its maximum over much of the support, most sharply when the CATE changes steeply over a narrow range---the threshold case in which equation~(\ref{eq:heterogeneity_bound}) is itself conservative (Section~\ref{sec:discussion}).

The right-hand side of equation~(\ref{eq:heterogeneity_bound}) involves $W_1$ in its original measurement units, multiplied by a clinical slope $L_{\text{clinical}}$; no further normalization of $W_1$ is required to obtain a clinically interpretable bound. We therefore adopt $W_1$ in its original units as the primary similarity metric throughout this paper. One could instead normalize $W_1$ by a within-effect-modifier scale measure such as the pooled interquartile range; we show in Section~\ref{sec:inference} that such normalization is mathematically redundant once equation~(\ref{eq:heterogeneity_bound}) is used for clinical calibration. The Lipschitz assumption underlying Proposition~\ref{prop:heterogeneity} is appropriate when the CATE function varies smoothly with the effect modifier; for threshold or step-function CATEs, $L_{\text{clinical}}$ can be arbitrarily large and the bound becomes conservative (Section~\ref{sec:discussion}).

\subsection{Estimation}\label{sec:estimation}

Given samples $\{X_{1,i}\}_{i=1}^{n_1}$ from region 1 and $\{X_{2,j}\}_{j=1}^{n_2}$ from region 2, the Wasserstein-1 distance is estimated using empirical distribution functions. Since $\hat{F}_1$ and $\hat{F}_2$ are piecewise constant (right-continuous step functions), the integral in equation~(\ref{eq:wasserstein}) reduces to a finite sum over the intervals between consecutive distinct values of the pooled sample:
\begin{equation}
\widehat{W}_1 = \sum_{k=1}^{m-1} |\hat{F}_1(x_{(k)}) - \hat{F}_2(x_{(k)})| \cdot (x_{(k+1)} - x_{(k)}),
\label{eq:estimator}
\end{equation}
where $\hat{F}_r$ denotes the empirical CDF (right-continuous) of region $r$ and $x_{(1)} < \cdots < x_{(m)}$ are the $m$ distinct values of the pooled sample (ties are pervasive when an effect modifier is recorded on a coarse grid, as with systolic blood pressure recorded in whole mmHg in Section~\ref{sec:application}). The estimator $\widehat{W}_1$ inherits the original measurement units of the effect modifier. The asymptotic distribution of $W_1$ is non-standard: del Barrio et al.\cite{delbarrio1999} showed that $\sqrt{n}\,W_1(\hat{F}_n, F) \xrightarrow{d} \int |B(F(x))|\,dx$, where $B$ is a standard Brownian bridge. Because the quantiles of this limiting distribution depend on the unknown $F$, no universal critical values are available and direct construction of asymptotic confidence intervals is impractical. We therefore employ the nonparametric percentile bootstrap with $B = 2{,}000$ replicates.\cite{efron1993} In the two-sample setting, when the set $\{x : F_1(x) = F_2(x) \in (0,1)\}$ has Lebesgue measure zero---for example, when the two CDFs cross only at isolated points---the Hadamard derivative of the $L_1$ functional at $F_1 - F_2$ is linear, and the ordinary nonparametric bootstrap is consistent by the functional delta method.\cite{fang2019} Where the CDFs coincide on an interval, including $F_1 = F_2$, the derivative is non-linear and the naive bootstrap is not consistent.\cite{dumbgen1993,fang2019,sommerfeld2018} Convergence occurs at rate $\sqrt{n_1 n_2/(n_1+n_2)}$ (see Appendix~\ref{app:asymptotic} for details). The estimator~(\ref{eq:estimator}) is computed in $O((n_1 + n_2) \log(n_1 + n_2))$ time, dominated by sorting the combined order statistics.

Because $W_1$ is a distance, $\widehat{W}_1$ is non-negative by construction and is typically positive even when the two regions are draws from the same distribution. We refer to the sampling distribution of $\widehat{W}_1$ under $F_1 = F_2$ as the null floor. It is obtained by drawing two independent resamples, of the two regions' actual sizes, from one region's empirical distribution---the anchor's in Section~\ref{sec:app_operability}---and recomputing $\widehat{W}_1$; because it depends on both sample sizes and on the source distribution, it is computed per region pair. An observed $\widehat{W}_1$ at or below its 95th percentile, the reference point used in Section~\ref{sec:app_operability}, is one the estimator would produce with appreciable probability from identical distributions, so it is not evidence of a difference; how large the difference could be is read from the upper confidence limit of $\widehat{W}_1$, equivalently of $\Delta_{\max}$.

\subsection{Interpretation and Clinical Calibration}\label{sec:inference}

A central feature of $W_1$ is its connection to treatment effect heterogeneity through Proposition~\ref{prop:heterogeneity}. This connection provides the basis for a clinically grounded interpretation framework rather than reliance on fixed numerical cutoffs.

We refer to the Lipschitz constant of the CATE function $\tau(x)$ as the clinical slope $L_{\text{clinical}}$. By definition, $L_{\text{clinical}}$ quantifies the maximum change in the treatment effect per unit change in the effect modifier, expressed on the outcome scale of the trial---for example, absolute 30-day mortality risk per year of age, or per mmHg of systolic blood pressure. Throughout this paper $L_{\text{clinical}}$ is quoted on the proportion scale of the outcome, so that $L_{\text{clinical}} = 1 \times 10^{-2}$ per year denotes one percentage point of absolute 30-day mortality per year of age; clinical margins $\Delta_{\text{clin}}$ stated in percentage points are divided by 100 before entering any calculation. When prior individual patient data (IPD) are available, meta-regression on treatment-by-covariate interaction estimates the average interaction slope.\cite{riley2010,fisher2017} An average slope is, in absolute value, a lower bound on the Lipschitz constant rather than the constant itself: for non-linear $\tau$, the local slope can exceed the fitted one. An $L_{\text{clinical}}$ specified from such an estimate should therefore carry a sponsor-specified multiplicative safety factor, in the spirit of margin discounting in non-inferiority design, so that the premise of Proposition~\ref{prop:heterogeneity} remains defensible. More commonly at the MRCT planning stage, published subgroup reports document the direction of interaction rather than a per-unit slope, and a precise numerical $L_{\text{clinical}}$ is unavailable. We address both situations below.

From Proposition~\ref{prop:heterogeneity}, the maximum potential difference in regional average treatment effects attributable to effect modifier distributional differences is
\begin{equation}
\Delta_{\max} = L_{\text{clinical}} \cdot W_1(F_1, F_2),
\label{eq:delta_max}
\end{equation}
expressed in the units of the trial outcome scale. The formula is invariant to within-effect-modifier scale normalization of $W_1$: any normalized form $W_1 / X$ multiplied by the same scale $X$ cancels in equation~(\ref{eq:delta_max}), so adopting $W_1$ in its original units yields the same clinical-scale calibration as any rescaled version. When $L_{\text{clinical}}$ is independently known, equation~(\ref{eq:delta_max}) provides a direct clinical-scale translation of the observed distributional difference. When $L_{\text{clinical}}$ is unknown, the relation can be inverted to ask which $L_{\text{clinical}}$ would be required for a given $W_1$ to produce a clinically meaningful $\Delta_{\max}$.

The estimator $\widehat{W}_1$ has a percentile bootstrap CI (Section~\ref{sec:estimation}), and inferential uncertainty propagates linearly to $\Delta_{\max}$ when $L_{\text{clinical}}$ is treated as fixed: $[\Delta_{\max,L}, \Delta_{\max,U}] = L_{\text{clinical}} \cdot [W_{1,L}, W_{1,U}]$. When $L_{\text{clinical}}$ itself is uncertain or unavailable, we solve equation~(\ref{eq:delta_max}) for $L_{\text{clinical}}$ at a pre-specified clinical margin $\Delta_{\text{clin}}$ (the minimal clinically important treatment effect difference, the non-inferiority margin, or the targeted overall effect), yielding the required clinical slope
\begin{equation}
L^* = \frac{\Delta_{\text{clin}}}{W_1(F_1, F_2)}.
\label{eq:lstar}
\end{equation}
$L^*$ is the smallest clinical slope under which the observed distributional difference could produce a clinically meaningful regional treatment effect difference. Because $L^*$ is a decreasing function of $W_1$, the bootstrap CI of $\widehat{W}_1$ gives the percentile interval $[\Delta_{\text{clin}}/W_{1,U},\, \Delta_{\text{clin}}/W_{1,L}]$ for $L^*$. If a plausible upper bound on $L_{\text{clinical}}$ from prior evidence (denoted $L_{\text{UB}}$) lies below $L^*$, the distributional difference cannot produce such a difference, provided $L_{\text{clinical}} \leq L_{\text{UB}}$; reading $L_{\text{UB}}$ against the lower limit $\Delta_{\text{clin}}/W_{1,U}$ rather than the point value carries the estimation uncertainty of $\widehat{W}_1$ into this comparison. The two pathways---forward calibration via equation~(\ref{eq:delta_max}) when $L_{\text{clinical}}$ is known and reverse comparison via equation~(\ref{eq:lstar}) against $L_{\text{UB}}$ when $L_{\text{clinical}}$ is unknown---cover the two evidence states encountered at the MRCT planning stage.

Specifying $L_{\text{clinical}}$ (or its upper bound $L_{\text{UB}}$) is a clinical judgment, informed but not determined by prior evidence, for which the sponsor carries responsibility---the same footing as a non-inferiority margin. Writing $R_a(x)$ for the risk under arm $a$, $\tau(x) = R_1(x) - R_0(x)$, so $|\tau'(x)| \leq |R_1'(x)| + |R_0'(x)|$: arm-level risk gradients constrain the modification slope only jointly, and a prognostic gradient alone guarantees no bound, prognostic and treatment-effect gradients being distinct quantities.\cite{ballman2015,kent2020} For the GUSTO-I case study presented in Section~\ref{sec:application}, the Fibrinolytic Therapy Trialists' meta-analysis\cite{ftt1994}, GUSTO-I subgroup analyses,\cite{gusto1993} and the GUSTO-I prognostic model\cite{lee1995} inform illustrative bounds of $L_{\text{UB,age}} \approx 1 \times 10^{-2}$ per year and $L_{\text{UB,SBP}} \approx 2 \times 10^{-3}$ per mmHg on the 30-day mortality scale; these exceed the empirical prognostic age gradient of approximately $6 \times 10^{-3}$ per year derived from the age-stratified mortality reported by Lee et al.,\cite{lee1995} though not the sum of two arm-level gradients of that size that the inequality above permits---a reference point rather than a guarantee. Conclusions drawn from the bounds are conditional on them; because $L^*$ is reported for every region pair, any alternative bound can be compared against it directly. Alternative specification strategies include the deft estimator of Fisher et al.\cite{fisher2017}

The pooling assessment proceeds by computing per-effect-modifier $W_1$ with bootstrap CIs across candidate partner regions, reading each distance against its null floor and its upper confidence limit (Section~\ref{sec:estimation}), and deriving $\Delta_{\max}$ when $L_{\text{clinical}}$ is known or $L^*$ at a pre-specified $\Delta_{\text{clin}}$ when $L_{\text{clinical}}$ is unknown. The framework is deliberately estimation-centered rather than testing-based: $W_1$ has no universal cutoff that separates ``similar'' from ``dissimilar,'' and what a given $W_1$ means clinically depends on the clinical slope $L_{\text{clinical}}$, which is external to the distributional comparison. By presenting per-pair $\Delta_{\max}$ (or $L^*$) with bootstrap CIs alongside the clinical context---the overall treatment effect size, the non-inferiority margin, the minimal clinically important difference---the framework provides quantitative inputs for sponsor and regulatory judgment. The judgment is not driven by a significance test.\cite{wasserstein2016}


% ---- Study 2 (estimation) design, from Section 3.2 of the manuscript (no results) ----
\subsection{Study 2: Estimation Properties}\label{sec:study1}

Study 1 compared the measures' ability to discriminate; Study 2 characterizes the estimator on which the calibration relies: bias, variability, and coverage probability of bootstrap confidence intervals for $\widehat{W}_1$. These properties underwrite the clinical calibration framework, as reliable estimation of $W_1$ directly determines the quality of $\Delta_{\max}$ estimates via equation~(\ref{eq:delta_max}). We report operating characteristics in the original units of the effect modifier, with each scenario specified at $\sigma = 10$; bias and CI width are therefore on the $W_1$ scale, while coverage is unitless and directly comparable across scenarios.

\subsubsection{Design}\label{sec:sim_design}

We designed seven systematic scenarios for methodological validation, examining controlled distributional differences in location, scale, scale combined with skewness (S6, a mean-matched log-normal), and a location--scale combination. Table~\ref{tab:scenarios} summarizes these scenarios with their true $W_1$ values, in the original units of the effect modifier.

\begin{table}[ht]
\centering
\caption{Study 2 simulation scenarios with true $W_1$ values.\label{tab:scenarios}}
\begin{tabular}{lllll}
\toprule
\textbf{ID} & \textbf{Description} & \textbf{Distribution 1} & \textbf{Distribution 2} & \textbf{True $W_1$} \\
\midrule
S1 & Null (identical) & $N(50, 10^2)$ & $N(50, 10^2)$ & 0.00 \\
S2 & Location 0.2$\sigma$ & $N(50, 10^2)$ & $N(52, 10^2)$ & 2.00 \\
S3 & Location 0.5$\sigma$ & $N(50, 10^2)$ & $N(55, 10^2)$ & 5.00 \\
S4 & Location 1.0$\sigma$ & $N(50, 10^2)$ & $N(60, 10^2)$ & 10.00 \\
S5 & Scale 1.5$\times$ & $N(50, 10^2)$ & $N(50, 15^2)$ & 3.99 \\
S6 & Skew (Log-normal) & $N(50, 10^2)$ & LogN$(3.787, 0.5^2)$ & 12.16 \\
S7 & Location + Scale & $N(50, 10^2)$ & $N(55, 15^2)$ & 5.83 \\
\bottomrule
\end{tabular}
\begin{tablenotes}
\item The Log-normal in S6 has $\sigma_{\log} = 0.5$ and $\mu_{\log} = \log 50 - \sigma_{\log}^2/2 = 3.787$, giving mean 50, CV $= 53.3\%$ and skewness $1.75$. S7 combines location and scale differences. True $W_1$ values for S1--S4 are exact ($W_1 = |\mu_1 - \mu_2|$ for equal-variance normals); S5 uses the closed form $W_1 = \sqrt{2/\pi}\,|\sigma_1 - \sigma_2|$; S6 and S7 are obtained by numerical quadrature of $W_1 = \int |F_1(t) - F_2(t)|\,\mathrm{d}t$, with the S6 integral split at $t = 0$ to accommodate the log-normal support.
\end{tablenotes}
\end{table}

For each scenario, we generated samples of size $n = 50$, 100, and 200 per region, reflecting sample sizes typical in MRCT regional subgroups. We performed 10{,}000 replications per scenario--sample size combination to ensure stable estimates of operating characteristics, with Monte Carlo standard errors below 0.005 for all reported proportions. Bootstrap confidence intervals were computed using $B = 2{,}000$ resamples. All simulations were conducted in R version 4.5.1.

Consistent with the estimation-centered framework, we evaluated the operating characteristics of the sample Wasserstein-1 estimator $\widehat{W}_1$ for the population value $W_1$, reported in the original units of the effect modifier:
\begin{enumerate}[1.]
\item Bias: $\text{Mean}(\widehat{W}_1) - W_1$
\item Root mean squared error (RMSE): $\sqrt{\text{Mean}[(\widehat{W}_1 - W_1)^2]}$
\item Coverage probability: Proportion of 95\% bootstrap CIs containing the true value
\item CI width: Mean width of bootstrap CIs, reflecting estimation precision
\end{enumerate}
